\Askhsh{16144}{2}
\begin{ylh}{1}
    \item 1.1. Η Έννοια του Διανύσματος
    \item 1.5. Εσωτερικό Γινόμενο Διανυσμάτων.
\end{ylh}
\begin{wrapfigure}[12]{r}{4cm} % 12 lines for the text to wrap around, 4cm width for the figure
\centering
\begin{tikzpicture}[scale=0.8] % Scale the figure to fit
    % Define points based on O=(0,0) and rhombus properties (side 4, angle A=60)
    % In a rhombus with angle A=60, triangle ABD is equilateral, so BD=4.
    % In right triangle AOD, AD=4, angle DAO=30 degrees.
    % AO = AD * cos(30) = 4 * (sqrt(3)/2) = 2*sqrt(3)
    % OD = AD * sin(30) = 4 * (1/2) = 2
    \tkzDefPoint(0, {2*sqrt(3)}){A} % A is on the positive y-axis
    \tkzDefPoint(-2, 0){B}           % B is on the negative x-axis
    \tkzDefPoint(0, {-2*sqrt(3)}){C} % C is on the negative y-axis
    \tkzDefPoint(2, 0){D}            % D is on the positive x-axis
    \tkzDefPoint(0,0){O}             % O is the origin

    % Draw the rhombus perimeter
    \tkzDrawPolygon[line width=0.8pt, darkgray](A,B,C,D)

    % Draw the diagonals
    \tkzDrawSegments[line width=0.8pt, darkgray](A,C B,D)

    % Draw points
    \tkzDrawPoints[size=4pt, fill=black](A,B,C,D,O)

    % Label points
    \tkzLabelPoint[above left](A){A}
    \tkzLabelPoint[left](B){B}
    \tkzLabelPoint[below right](C){$\Gamma$}
    \tkzLabelPoint[right](D){$\Delta$}
    \tkzLabelPoint[above right=1pt](O){O} % Adjust position for O label

    % Mark the angle A = 60 degrees (angle DAB)
    \tkzMarkAngle[arc=ll, size=0.6cm, fill=maincolor!10, draw=maincolor](D,A,B)
    \tkzLabelAngle[pos=0.9, secondarycolor](D,A,B){$60^\circ$}
\end{tikzpicture}
\end{wrapfigure}
Δίνεται ρόμβος ΑΒΓΔ με κέντρο Ο, πλευρά $4$ και $\hat{A} = 60^\circ$. Να υπολογίσετε τα εσωτερικά γινόμενα :
\begin{alist}
\item $\vec{AB} \cdot \vec{A\Delta}$
\item $\vec{A\Delta} \cdot \vec{B\Gamma}$
\item $\vec{O\Delta} \cdot \vec{AO}$
\item $\vec{O\Delta} \cdot \vec{OB}$
\item $\vec{A\Delta} \cdot \vec{\Gamma\Delta}$ \monades{25}
\end{alist}